\chapter{Conclusiones}\label{cap:conclusiones}

\section{Consecución de Objetivos}
Tras la finalización del desarrollo de TraDeck, se puede afirmar que se han alcanzado satisfactoriamente los objetivos planteados en la fase de diseño:

\begin{itemize}
	\item \textbf{Mercado P2P funcional}: Se ha logrado implementar una plataforma donde los usuarios pueden interactuar directamente para el intercambio de activos, las cartas NFT, sin intermediarios centrales.
	\item \textbf{Seguridad y Custodia}: El uso de contratos inteligentes con protocolos \textit{Escrow} y \textit{Atomic Swaps} garantiza que las transacciones sean seguras y atómicas, protegiendo la integridad de los activos.
	\item \textbf{Persistencia Híbrida Eficiente}: Se ha validado la arquitectura que combina el \textit{ledger} de Ethereum para la titularidad y la red IPFS para los datos multimedia, optimizando drásticamente los costes de almacenamiento.
	\item \textbf{Accesibilidad Web3}: La interfaz desarrollada permite una interacción fluida con la blockchain mediante MetaMask, abstrayendo la complejidad técnica para el usuario final.
\end{itemize}

\section{Limitaciones del Sistema}
A pesar del éxito del prototipo, se han identificado limitaciones críticas que definen el estado actual de la tecnología:

\begin{itemize}
	\item \textbf{Latencia en la Recuperación de Datos Históricos}: Durante la implementación, se detectó un cuello de botella al consultar el estado de las cartas mediante eventos. Inicialmente, la búsqueda era progresiva, lo cual funcionó genial el mismo día, al probarlo, pero a medida que los activos quedaban "hundidos" miles de bloques atrás en la cadena, la respuesta del frontend se volvía inaceptable, hasta 720 llamadas a la cadena hicieron falta para conseguir una sola carta. Esto obligó a una refactorización hacia el uso de estructuras de datos optimizadas (como el patrón DTO \textit{CardData}, \textit{struct} en los contratos) para realizar consultas directas al estado del contrato en una sola petición RPC.
	\item \textbf{Dependencia de Proveedores}: Aunque el sistema es descentralizado, la experiencia de usuario actual depende de la disponibilidad de servicios de alojamiento gratuitos para el proxy y el frontend, los cuales pueden entrar en modo reposo.
	\item \textbf{Costes de Gas}: En la red principal, el coste de las transacciones de escritura sigue siendo un factor limitante para operaciones de bajo valor.
\end{itemize}

\section{Trabajo Futuro}
Para evolucionar TraDeck hacia un producto comercial, se proponen las siguientes líneas de mejora:

\begin{itemize}
	\item \textbf{Migración a Capa 2 (L2)}: Implementar el sistema en redes como Optimism o Arbitrum para reducir los costes de gas en más de un 90\%. Es decir, saltar de una testnet como Sepolia a una red real.
	\item \textbf{Gobernanza mediante DAO (Organizaciones Autónomas Descentralizadas)}: Permitir que la comunidad decida sobre la creación de nuevas colecciones o modificaciones en las reglas del mercado.
	\item \textbf{Integración de Oráculos}: Utilizar servicios como Chainlink para verificar el valor de mercado externo de las cartas y ofrecer estadísticas en tiempo real.
	\item \textbf{Indexación Avanzada}: Implementar \textit{The Graph} para solucionar definitivamente los problemas de latencia en consultas históricas complejas. Fue una idea que nosotros mismos planteamos para el primer prototipo pero descartamos porque subiría la complejidad del proyecto en muchas horas.
	\item \textbf{Monetización real}: Pasar de nuestra propia moneda, TDC, a usar criptomonedas reales, permitiendo compra venta directamente por la plataforma sin necesidad de servicios externos.
\end{itemize}

\section{Valoración Personal}
El desarrollo de este proyecto ha supuesto un desafío técnico importante que ha superado las expectativas iniciales. La carga de trabajo real ha excedido las 60 horas estimadas debido a la curva de aprendizaje de la Web3 y, especialmente, a la resolución de problemas imprevistos de rendimiento en la comunicación con la blockchain. 

La necesidad de refactorizar la lógica de obtención de datos del frontend fue la tarea que más tiempo consumió, pero también la que mayor aprendizaje aportó sobre la arquitectura de sistemas distribuidos. El sistema ya estaba funcionando, todo iba bien, y no fue hasta el día siguiente que empezaron a salir errores, ya que se necesitaba que las cartas se hundieran en la cadena y eso no depende de nosotros si no del propio paso del tiempo. Por suerte pudimos detectar el error, conseguimos dar con la solución usando un patrón DTO y una llamada directa al estado y no a eventos históricos.

En conclusión, TraDeck no solo es una plataforma de intercambio, sino una prueba de concepto sólida sobre cómo la inmutabilidad y la descentralización pueden transformar sectores tradicionales como el coleccionismo, haciendo más fiable y veraz el intercambio de activos entre particulares.